\section{2026 年 11 月核验增补}

本节按论文首次提交时间排列。每项方法附有可核验的外部入选依据。

% METHOD: 2510.08566
\paragraph{D²GS：分别约束少视图重建的近端过拟合与远端欠拟合}
2025.10~\cite{2510_08566}。少视图训练时，近相机的高密度区域可能容纳过多高斯，远端反而缺少几何。D²GS 在 DropGaussian 骨干上用深度与局部原语密度控制近处分布的随机失活概率，再用单目深度给远端加带遮罩的图像监督。测试还用跨重复训练的原语分布一致性辅助观察稳定程度。

论文在 LLFF 和 Mip-NeRF 360 及补充的多视图条件比较，部分 DropGaussian 数值来自作者自己的重训。固定深度阈值与权重难以适配所有场景，跨模型一致性也不等于动态新视图感知稳定。原文的机构署名包括 Insta360 Research。

\textbf{入选依据} 论文署名单位包括 Insta360（\href{https://arxiv.org/html/2510.08566}{论文原文}）。

% METHOD: 2605.30342
\paragraph{GAVIS：用方向相关可见性评价下一次主动观察}
2026.05~\cite{2605_30342}。GAVIS 将已训练的高斯场和相机方向信息转成各向异性可见性场，用球谐解析计算未充分观察的方向。规划器根据候选视点的不确定性选择观察，无需为每次规划重新训练一个可见性网络。

论文以统一的候选采样和规划管线对比不确定性模块，在 NeRF Synthetic、太空模型、HM3D 和 Gibson 场景验证，并加入 FisherRF 与 VIMC 作后置校准。比较依赖仿真中可行候选、已有视图与可见性定义，不能把其速度直接等同于实物机器人闭环周期。公开实现包含核心模块和运行脚本。

\textbf{入选依据} \href{https://github.com/xsj01/gavis}{公开实现}已找到与论文对应的代码文件。

% METHOD: 2607.05765
\paragraph{Image2Sim：把被动图像与视频变成可交互的导航训练环境}
2026.07~\cite{2607_05765}。Image2Sim 从多视图观察前馈重建带特征的高斯，结合一步生成的全景 RGB-D、碰撞约束运动及自动指令构造，让智能体能在场景中逐步观察与行动。由此可以生成用于视觉语言导航训练的轨迹，不需要逐个场景人工建模。

论文构建 19936 个不含评测或测试环境的训练场景，并分别测新视图渲染质量和导航。作者在真实 Stretch 3 上仅以有限次数评估两类指令，并在附录说明主要导航增益来自合成数据的规模和多样性。作者提供生成、训练和评估相关代码。

\textbf{入选依据} \href{https://github.com/mrzihan/image2sim}{公开实现}已找到与论文对应的代码文件。

% METHOD: 2607.17965
\paragraph{探索式增密：缓解高斯重建中的远侧与近侧模糊陷阱}
2026.07~\cite{2607_17965}。该工作把部分模糊归因于视线方向有效位置梯度不足，以及遮挡导致的透明度混合梯度衰减。随机播种补充沿深度方向的位置候选，随机分裂让部分原语不必等到梯度超过阈值才增密，分别针对两类问题。

论文以多个基准比较两种方式及其组合，消融显示它们在遮挡和远侧区域互补。改善仍依据所测场景和训练预算，随机探索并不保证每个新增原语都落到真实表面。作者公开包含训练、渲染和评估入口的实现。

\textbf{入选依据} \href{https://github.com/chengbo-wang/exploration-for-gs}{公开实现}已找到与论文对应的代码文件。

% METHOD: 2608.01178
\paragraph{DynActiveGS：分别衡量结构不确定性和动态干扰}
2026.08~\cite{2608_01178}。DynActiveGS 从动态环境的 RGB-D 观测估计像素不确定性，控制观测在高斯重建中的权重，再分别用结构和运动不确定性在 Voronoi 图上选视点。规划既要补足静态几何的覆盖，也要考虑动态物体可能造成的遮挡。

论文在 Social-MP3D 与动态增强的 Gibson 模拟场景检查重建与保留视图，保留视点不出现在各方法的训练轨迹里。结果对应模拟动态环境，传感器噪声和未建模动态仍需实机验证。

\textbf{入选依据} 论文署名单位包括 Huawei（\href{https://arxiv.org/abs/2608.01178}{论文原文}）。

% METHOD: 2608.20602
\paragraph{通过同步稀疏光场为随手采集的高斯重建补充角度信息}
2026.08~\cite{2608_20602}。手机多摄、双目和光场设备能在同一曝光获得相邻视角。作者把这种同步的角度信息保留在高斯重建中，并制作子透镜影像可以重叠的复用光场原型，试图兼顾单视图空间分辨率和角度采样。同步视角对动态场景尤其有用，因为移动单相机不能回到同一时刻补采。

论文在匹配曝光数和改动角间距的合成实验上比较图像与深度，并用真实静态、动态和原型采集做验证。Adobe 署名可核实，仓库列有数据、训练和运行入口。部分动态流程依赖另行申请的非商用模型权重，不能把公开代码理解为所有依赖都无条件开放。

\textbf{入选依据} 论文署名单位包括 Adobe（\href{https://arxiv.org/html/2608.20602v1}{论文原文}）。\href{https://github.com/shamus-li/lightfield-gaussian-splatting}{公开实现}已找到与论文对应的代码文件。

% METHOD: 2609.22941
\paragraph{D3GS：用米制深度、跨视图特征和扩散辅助少视图优化}
2026.09~\cite{2609_22941}。D3GS 首先把来自 MapAnything 与束调整的稀疏米制深度补全后初始化高斯，再用深度、DINO 特征与扩散修复为优化提供互补监督。几何先验、跨视图对齐和生成外观分别约束不同误差来源。

论文报告少视图设定下的几何和画质比较。若初始米制深度失准，后续初始化可能继承误差；扩散生成补充的图像外观也需要通过跨视图观测判断是否对应真实几何。

\textbf{入选依据} 论文署名单位包括 ByteDance（\href{https://arxiv.org/abs/2609.22941}{论文原文}）。

% METHOD: 2610.07110
\paragraph{MoonGS：从双视图和语义先验重建低纹理月面}
2026.10~\cite{2610_07110}。MoonGS 从一对带相机参数和语义图像的月面视图直接预测高斯。方法利用视觉基础模型的深度特征确定基元位置，在非重叠或缺少纹理的区域以语义先验改善深度估计，并用深度排序约束稳定背景。启发式视图重采样用于补足稀疏观测。

作者在 LuSNAR 和合成 MoonBlender 上以重训的前馈新视图模型为基线，LuSNAR 的 PSNR 由所比基线的最高 18.89 dB 提高到 23.80 dB。真实嫦娥图像只用于定性比较。该增益不能直接解读为真实月面的米制几何误差下降，两幅输入也限制了大面积遮挡区的证据。

\textbf{入选依据} \href{https://github.com/inrobots/moonblender}{公开实现}已找到与论文对应的代码文件。
